\chapter{Spaces, Vectors, and Inner Products}
\label{ch:spaces-and-inner-products}

A vector $x \in \R^d$ is not merely a container. Its dimension bounds capacity; an inner product defines similarity; a basis decides which variations count as independent. The geometry of that space is already a theory of what can be compared.

Raw observations---pixels, tokens, sensor readings---become usable only after they are placed in such a space. Figure~\ref{fig:inner-product} is the geometry that later chapters reuse as a score.

\begin{figure}[htbp]
  \centering
  \begin{tikzpicture}[scale=1.1]
  \draw[->] (0,0) -- (3.2,0) node[right] {$x$};
  \draw[->] (0,0) -- (0,2.6) node[above] {$y$};
  \draw[thick,accent] (0,0) -- (2.4,1.5) node[above right] {$u$};
  \draw[thick] (0,0) -- (2.1,0) node[below] {$v$};
  \draw[dotted] (2.4,1.5) -- (2.1,0);
\end{tikzpicture}

  \moicaption{Inner product in the plane}{The vector $u$ and a reference direction $v$. The dotted segment is the drop that makes $\langle u,v\rangle$ visible as a signed length along $v$.}
  \label{fig:inner-product}
\end{figure}

\begin{definition}[Inner product]
\label{def:inner-product}
An \emph{inner product} on a real vector space $V$ is a bilinear map
$\langle \cdot,\cdot\rangle \colon V\times V\to\R$ that is symmetric and
positive definite: $\langle x,x\rangle\ge 0$, with equality if and only if $x=0$.
The induced norm is $\lVert x\rVert=\sqrt{\langle x,x\rangle}$.
\end{definition}

\begin{theorem}[Cauchy--Schwarz]
\label{thm:cauchy-schwarz}
For all $u,v\in V$,
\[
\lvert\langle u,v\rangle\rvert \le \lVert u\rVert\,\lVert v\rVert.
\]
\end{theorem}

\begin{proof}
If $v=0$ the claim is immediate. Otherwise set
$t=\langle u,v\rangle/\lVert v\rVert^2$ and expand the inequality
$\langle u-tv,u-tv\rangle\ge 0$.
\end{proof}

\begin{corollary}[Cosine bound]
\label{cor:cosine-bound}
If $u$ and $v$ are nonzero, the cosine of the angle between them satisfies
$\lvert\cos\theta\rvert\le 1$, where
\[
\cos\theta=\frac{\langle u,v\rangle}{\lVert u\rVert\,\lVert v\rVert}.
\]
\end{corollary}

\begin{example}[Cosine similarity]
\label{ex:cosine-similarity}
Token or image embeddings are compared by the same cosine. A score near $1$
means the two vectors point the same way---not that they are close in Euclidean
distance.
\end{example}

\begin{remark}
Later chapters return to Definition~\ref{def:inner-product} when attention,
retrieval, and contrastive learning reuse inner products as scores.
\end{remark}
